\documentclass[11pt]{article}

% ---------- Packages ----------
\usepackage[utf8]{inputenc}
\usepackage[T1]{fontenc}
\usepackage{lmodern}
\usepackage[margin=1in]{geometry}
\usepackage{amsmath,amssymb}
\usepackage{siunitx}
\usepackage{graphicx}
\usepackage{float}
\usepackage{caption}
\usepackage{subcaption}
\usepackage{booktabs}
\usepackage{microtype}
\usepackage{fancyhdr}
\usepackage{xcolor}
\usepackage[colorlinks=true,linkcolor=blue!55!black,citecolor=blue!55!black,urlcolor=blue!55!black]{hyperref}

% ---------- Spacing / look ----------
\setlength{\parskip}{0.55em}
\setlength{\parindent}{0pt}
\captionsetup{font=small,labelfont=bf}
\graphicspath{{figures/}}

% ---------- Header / footer ----------
\pagestyle{fancy}
\fancyhf{}
\lhead{\small EE132 --- Automatic Control}
\rhead{\small Lab 7}
\cfoot{\thepage}
\renewcommand{\headrulewidth}{0.4pt}

% siunitx setup
\sisetup{detect-all,per-mode=symbol}

% ---------- Title ----------
\title{\vspace{-1.0cm}\textbf{Lab 7 Report}\\[0.3em]
\large Proportional--Velocity (PV) Position Control of the Quanser QUBE-Servo}
\author{Kaushik Vada \and Andrew Kong \and Aren Muniz \and Daniel Lopez \and Kai Ong}
\date{May 29, 2026}

\begin{document}
\maketitle
\thispagestyle{fancy}

% ---------- Abstract ----------
\begin{abstract}
\noindent This lab focuses on using a proportional--velocity (PV) controller for QUBE-Servo
position control. A Simulink model was developed and tested to evaluate the effects of
proportional and derivative gains on system performance. Controller gains were designed
using second-order system specifications and implemented on the servo. The resulting
position response was analyzed to verify that the desired peak time, overshoot, and overall
stability were achieved.
\end{abstract}

% =====================================================================
\section{System Built}
% =====================================================================
The system built in this lab was a proportional--velocity (PV) position control system for
the Quanser QUBE-Servo. The system was implemented in Simulink and consisted of a reference
signal generator, a proportional controller, a velocity feedback path with a high-pass
filter, and the QUBE-Servo plant. A square-wave position command was used as the reference
input, while the measured servo position was fed back to the controller. The proportional
gain was used to improve tracking performance, while the derivative gain provided damping
and reduced overshoot. The completed system was tested to evaluate position-control
performance and to verify the controller gains designed from second-order system
specifications.

\begin{figure}[H]
    \centering
    \includegraphics[width=0.95\textwidth]{model.jpg}
    \caption{QUARC/Simulink PV position control model for the QUBE-Servo. The proportional
    gain $K_p$ acts on the position error while the velocity feedback path uses a high-pass
    filter ($15s/(s+150)$) and derivative gain $K_d$ for damping; the controller output
    drives the motor through the HIL Write Analog block.}
    \label{fig:model}
\end{figure}

% =====================================================================
\section{In-Lab Exercises and Results}
% =====================================================================

\subsection{Effect of the Proportional and Derivative Gains}
The motor voltage command and the position-tracking response were observed while sweeping
the proportional gain $K_p$ and the derivative gain $K_d$. Figure~\ref{fig:voltage} shows
the motor voltage for the square-wave command, and Figure~\ref{fig:position} shows the
reference and measured position together.

\begin{figure}[H]
    \centering
    \includegraphics[width=0.82\textwidth]{voltage.jpg}
    \caption{Motor voltage command $V_m$ for the square-wave position reference.}
    \label{fig:voltage}
\end{figure}

\begin{figure}[H]
    \centering
    \includegraphics[width=0.82\textwidth]{position.jpg}
    \caption{Position-tracking response: the square-wave reference (yellow) and the measured
    servo position (blue) following it.}
    \label{fig:position}
\end{figure}

Sweeping the proportional gain produced the following observations:
\begin{itemize}
    \item $K_p = 1$: the amplitude is not as high --- the previous response was around
    $2.5$, while this is around $1$.
    \item $K_p = 2$: the output has more oscillations; at least, the smaller oscillations
    are now larger.
    \item $K_p = 3$: even more oscillations than when $K_p = 2$.
    \item $K_p = 4$: even more oscillations than when $K_p = 3$.
\end{itemize}
As $K_p$ increases there are more oscillations, which shows that the servo is moving more
frequently the higher the value of $K_p$. In contrast, as $K_d$ increases, the output
result becomes more similar to the reference.

\subsection{Controller Gain Design}
Matching the QUBE-Servo closed-loop transfer function to the standard second-order form
yields the proportional and derivative gains
\begin{equation}
    K_p = \frac{\tau\,\omega_n^{2}}{K},
    \qquad
    K_d = \frac{2\zeta\omega_n\tau - 1}{K}.
    \label{eq:gains}
\end{equation}
Using the model parameters $\omega_n = 32.3$, $\zeta = 0.76$, $K = 23.2$, and
$\tau = 0.13$,
\begin{align}
    K_p &= \frac{(0.13)(32.3)^{2}}{23.2} = 5.85, \\[4pt]
    K_d &= \frac{2(0.76)(32.3)(0.13) - 1}{23.2} = 0.232 .
\end{align}
Therefore the designed gains are
\begin{equation}
    K_p = \SI{5.85}{\volt\per\radian}
    \qquad\text{and}\qquad
    K_d = \SI{0.232}{\volt\per(\radian\per\second)} .
\end{equation}

% =====================================================================
\section{Analysis}
% =====================================================================
\textbf{Results analysis.} The PV controller successfully controlled the position of the
QUBE-Servo and demonstrated the effects of proportional and derivative feedback on system
performance. Increasing the proportional gain improved the speed of the response and reduced
tracking error, but excessive proportional gain increased overshoot and oscillatory
behavior. Increasing the derivative gain added damping to the system, reducing overshoot and
improving stability. Using the calculated gains, the servo response closely followed the
reference position with a fast rise time and minimal overshoot. Overall, the experimental
results were consistent with the theoretical second-order system model and demonstrated that
PV control can effectively improve transient response while maintaining stable operation.

\textbf{Issues encountered.} None encountered.

% =====================================================================
\section{Conclusion}
% =====================================================================
In this lab, a proportional--velocity (PV) controller was designed and implemented to
control the position of the QUBE-Servo system. The effects of proportional and derivative
gains on the transient response were investigated, and controller gains were calculated
using second-order system design specifications. Experimental testing showed that
proportional gain increased response speed while derivative gain improved damping and
reduced overshoot. The final controller achieved stable position tracking and demonstrated
the effectiveness of PV control for servo position regulation. Overall, the lab provided
practical experience in control system design, tuning, and performance analysis using a real
electromechanical system.

\end{document}
